\fullbleedpage{assets/cover-third-edition.png}

\frontmatter
\chapter*{第二版前言}

去年春天，一些问题的提出与对话，促使我们动笔编写了第一版《北京大学工学院新生生存手册》。它源于工学院学生会对新生入学遇到的信息壁垒、沟通效率低下、成长陷入迷茫等问题的关怀。第一版完成后，我们收到了许多反馈与建议，同时也意识到：一份真正有价值的手册，不应止步于初稿，它应当与时俱进，精进自身，回应变化。

\vspace{10pt}

所以，我们决定继续写下去。

\vspace{10pt}

这一版，我们保留了实用信息的基础，也在内容、结构与表达方式上进行了调整和补充。我们尽己所能，尝试用更细致的语言，回应你可能遇到的问题：从上课到科研，从宿舍到教室，从学业规划到情绪照料。我们希望它不仅仅是一本指南，更是一份陪伴你适应、思考、探索的同行记录。

\vspace{10pt}

笔者真诚地希望，在你打开这本手册的瞬间，能感受到属于每一位北大工院人的温度。大学不只是课程、成绩，不只是科研、竞赛。它也关于在五四体育场迎着朝霞的晨跑、在未名湖边夜深的对话、在邱德拔体育馆挥洒的汗水、在百周年纪念讲堂的舞台上绽放的光芒。成长的模样本就多种多样，而你会在自己的节奏中，找到最真实的路径。

\vspace{10pt}

这本手册写给你，也写给那个一年前的我们。愿它能为你照亮某个时刻，哪怕只是一小段路。

\vspace{10pt}

愿你在燕园里，心有所归，目有繁星。

\begin{flushright}
	郭濠源
	
	2025年8月
\end{flushright}

\chapter*{前言}
那是一个寒冷的初春。一条不起眼树洞上的几条回复几乎可以说直接导致了这份手册的产生。2024年3月，工学院学生会学术部开始对本手册进行筹划。经过一个暑期的编写，新生版的手册终于能够呈现在初入燕园的同学们手中。

\vspace{10pt}

笔者首先感谢学术部倪昊、武昱达同学抽出暑假时间来编写这份手册，感谢学院内各位老师的悉心指导和谢谊锋、吴秉宪以及手册中约稿内容的作者们等诸多同学的大力支持；没有他们的努力和智慧，这份手册不会如此快速、高效地呈现在大家面前。笔者也要感谢入校以来与自己沟通、交流过手册编写的所有人，与他们的交流为手册提供了灵感和准备。在这里，我们同样期待着每一位愿意向大家传递信息、分享经验，让本手册更好、更充实的各位加入我们。更新迭代是手册的必经之路，只有这样，它才能跟上时代的发展，才能更好地为更多同学的成长与发展服务。

\vspace{10pt}

回到最初，是什么引发了我们编写本手册的想法？我回顾在工学院度过的这两年，深感学院在经历剧烈的变化，从招生人数的剧增、培养方案的改写、教师队伍的“握手计划”、新奥大楼的建成等方面均可见一斑。招生人数的增加本应该让同学之间的相互交流学习变得更加轻松，但我看到了很多同学没有办法从现有的渠道中获得信息，也不愿意找老师、身边的同学进行交流，最终导致错过了很多机会；我也看到了有些同学之间的陌生、小范围抱团却不免闭塞的交流圈子；我还注意到老师们、学长学姐们的分享往往是点对点的，会消耗大量不必要的时间精力，且只能覆盖问题在个体上的投射，需要多个角度才能知晓全貌……我始终认为：从我入学直至现在两年的时间，工学院同学们在信息获取和交流方面的低效状况未曾有大的改变，却亟待解决。本手册希望做出一些力所能及的贡献。

\vspace{10pt}

当然手册的任务也不仅如此。手册还希望帮助大家在本科四年里寻求一些问题的答案：我是怎样的人？我喜欢什么？我想要去做些什么？新的时代快速发展的浪潮下，这些问题是需要想清的，因此我们可以尽早开始思考。

\vspace{10pt}

最后笔者想向大家明确一点：本手册不是《高分保研奖学金捷径手册》或者《内卷宝典》。本手册希望帮助大家的，绝不是仅仅拿一个高高的GPA（Grade Point Average，平均学分绩点）、发SCI（Science Citation Index，科学引文索引）论文、收到顶尖学校的直博录取诸如此类能被量化的指标。工学院的课程压力不小，但生活不能只有专业课和成绩，燕园里的其他风景也同样精彩。我希望这份手册能够带给大家一点恰当的提示、一些新鲜的思考，让同学们能看到人生道路上不同的风景。至少在读过本手册之后，各位应该能够有勇气和智慧，去发现并挑战一些比既定指标更重要的追求。
\begin{flushright}
	赵丹枫
	
	2024年8月
\end{flushright}

\chapter*{前言}
开辟鸿蒙，谁为学种？都只为工学情浓。趁着这迎新天、报到日、欢乐时，试着剧透。因此上，撰写这震古烁今的《生存册》。——《工学梦曲》

\vspace{10pt}

欢迎各位新同学进入北京大学工学院！

\vspace{10pt}

《北京大学工学院新生生存手册》（以下简称“手册”）的撰写是一个长期项目，由工学院学生会全权负责。考虑到手册编写的目的是尽可能及时造福同学们，我们决定，1.0版本的手册面向的群体是本科新生，只会撰写最核心、最基础的一些内容，加之一些新生需要知道的事宜，以帮助刚刚入学的各位coeer快速弥合信息差，拓宽知识面，感受工院文化，找到属于自己的生活和学习节奏。

\vspace{10pt}

本手册旨在汇集学院老师、学长学姐们的经验和智慧，以一种自由且简洁的方式，为大家呈现直接、真实、系统性的工学院学习生活重要信息，减少大家搜索信息、辨别信息、整理信息的时间消耗，把更多的时间精力投入到对自己生涯的思考和规划、具体知识的学习和有意义的大学活动中，收获一段真正充实而精彩的大学生涯。当然，我们相信大家想要了解的信息不止于手册中包含的内容，因此我们在手册最后提供了问卷二维码，欢迎同学们进一步提问，或对我们的工作提出更好的建议。

\vspace{10pt}

手册自2024年3月开始筹备编写，如今尚处于“幼年”。希望同学们能从手册中获得一些启发，也希望更多的同学积极加入工学院学生会，为手册的编写贡献自己的智慧和力量！
\begin{flushright}
	倪昊
	
	2024年8月
\end{flushright}

\chapter*{声明}

《北京大学工学院新生生存手册》是一份由学生自发策划、组织编写的手册，版权属于北京大学工学院学生会。未经北京大学工学院学生会书面许可，任何组织或个人不得违反相应版权条例抄袭、转载、摘编、修改本书内容；不得将本书用于商业目的；不得对本手册原意进行曲解、修改和未授权的大范围分发。对于未经授权传播手册全部或部分内容而造成的各种问题，手册编者概不负责。

\vspace{10pt}

该手册集合众多学长学姐的经验和观点，并尽可能将各种观点统一在手册的框架下。这样做的目的，是希望向读者传达更多可供参考的观点和意见。请读者注意：本手册不构成任何明确的行动建议，不保证手册中的信息和方法始终正确、有效。编者不承担由手册及其内容产生的衍生责任。
\begin{flushright}
	北京大学工学院学生会
	
	2026年8月
\end{flushright}

\tableofcontents
